% 型5：高速ドリル型
% 短い問いを多数並べ，答えを段末の帯にまとめる．
% 計算の速さと正確さを鍛える用途に向く．解説は付けない．
\documentclass[lualatex,paper=b5j,fontsize=10pt,twoside]{jlreq}

\usepackage{surikumi-mondaishu}

\MathMagazineSetup{
  feature={},
  series={ドリル／数学 II},
  series-position=left,
  pattern=weave,
  title={微分の計算},
  author={入試基礎講座},
  author-font=mincho,
  author-position=right,
  title-fill=black,
  deck={1 問 15 秒を目安に解く．答は各段の下にまとめてあるので，手で隠して進めること．}
}

\begin{document}

\MakeMathMagazineTitle

\MMSubheading{A\quad 多項式の微分}

次の関数を $x$ について微分せよ．

\begin{msdrill}
  \item $y=x^3$
  \item $y=5x^2$
  \item $y=x^3-2x$
  \item $y=4x^3+3x^2-x+7$
  \item $y=(x+1)^2$
  \item $y=(2x-1)^3$
  \item $y=x^2(x-3)$
  \item $y=\dfrac{x^4}{2}-\dfrac{x^2}{4}$
\end{msdrill}

\begin{msanswerstrip}
  \MSAns{1}{$3x^2$}
  \MSAns{2}{$10x$}
  \MSAns{3}{$3x^2-2$}
  \MSAns{4}{$12x^2+6x-1$}
  \MSAns{5}{$2x+2$}
  \MSAns{6}{$6(2x-1)^2$}
  \MSAns{7}{$3x^2-6x$}
  \MSAns{8}{$2x^3-\dfrac{x}{2}$}
\end{msanswerstrip}

\MMSubheading{B\quad 接線}

曲線上の指定された点における接線の方程式を求めよ．

\begin{msdrill}
  \item $y=x^2$，点 $(1,1)$
  \item $y=x^2$，点 $(-2,4)$
  \item $y=x^3$，点 $(1,1)$
  \item $y=x^3-3x$，点 $(0,0)$
  \item $y=x^2-2x$，点 $(3,3)$
  \item $y=-x^2+4$，点 $(1,3)$
\end{msdrill}

\begin{msanswerstrip}
  \MSAns{1}{$y=2x-1$}
  \MSAns{2}{$y=-4x-4$}
  \MSAns{3}{$y=3x-2$}
  \MSAns{4}{$y=-3x$}
  \MSAns{5}{$y=4x-9$}
  \MSAns{6}{$y=-2x+5$}
\end{msanswerstrip}

\MMSubheading{C\quad 増減}

次の関数が極値をとる $x$ の値をすべて求めよ．極値をとらない場合は
「なし」と答えよ．

\begin{msdrill}
  \item $y=x^2-6x$
  \item $y=x^3-3x$
  \item $y=x^3+3x$
  \item $y=x^3-6x^2+9x$
  \item $y=x^4-2x^2$
  \item $y=x^3$
\end{msdrill}

\begin{msanswerstrip}
  \MSAns{1}{$x=3$}
  \MSAns{2}{$x=\pm 1$}
  \MSAns{3}{なし}
  \MSAns{4}{$x=1,\,3$}
  \MSAns{5}{$x=0,\,\pm 1$}
  \MSAns{6}{なし}
\end{msanswerstrip}

\MMNote{C の 3 と 6 は，$y'\geq 0$ が常に成り立つため増減が変わらない．$y'=0$ となる $x$ があっても極値とは限らない．}

\MMSubheading{D\quad 定積分}

次の定積分の値を求めよ．

\begin{msdrill}
  \item $\displaystyle\int_{0}^{1}x^2\,dx$
  \item $\displaystyle\int_{0}^{2}(3x^2-2x)\,dx$
  \item $\displaystyle\int_{-1}^{1}x^3\,dx$
  \item $\displaystyle\int_{-1}^{1}(x^2+1)\,dx$
  \item $\displaystyle\int_{1}^{3}(x-1)(x-3)\,dx$
  \item $\displaystyle\int_{0}^{1}(2x+1)^2\,dx$
\end{msdrill}

\begin{msanswerstrip}
  \MSAns{1}{$\dfrac{1}{3}$}
  \MSAns{2}{$4$}
  \MSAns{3}{$0$}
  \MSAns{4}{$\dfrac{8}{3}$}
  \MSAns{5}{$-\dfrac{4}{3}$}
  \MSAns{6}{$\dfrac{13}{3}$}
\end{msanswerstrip}

\MMNote{D の 3 は奇関数の対称区間上の積分であり，計算せずに $0$ と分かる．5 は $\dfrac{(\beta-\alpha)^3}{6}$ の公式が使える．}

\MMSubheading{E\quad 面積}

次の曲線と直線で囲まれた部分の面積を求めよ．

\begin{msdrill}
  \item $y=x^2$ と $y=x$
  \item $y=x^2-4$ と $x$ 軸
  \item $y=x^2$ と $y=2x+3$
  \item $y=x^3-x$ と $x$ 軸（$0\leq x\leq 1$ の部分）
\end{msdrill}

\begin{msanswerstrip}
  \MSAns{1}{$\dfrac{1}{6}$}
  \MSAns{2}{$\dfrac{32}{3}$}
  \MSAns{3}{$\dfrac{32}{3}$}
  \MSAns{4}{$\dfrac{1}{4}$}
\end{msanswerstrip}

\end{document}
